\documentclass[11pt,a4paper]{article}

\usepackage[T1]{fontenc}
\usepackage[utf8]{inputenc}
\usepackage{lmodern}
\usepackage{geometry}
\usepackage{setspace}
\usepackage{booktabs}
\usepackage{longtable}
\usepackage{array}
\usepackage{graphicx}
\usepackage{xcolor}
\usepackage{hyperref}
\usepackage{enumitem}
\usepackage{float}
\usepackage{amsmath}
\usepackage{caption}

\geometry{margin=2.5cm}
\setstretch{1.15}
\hypersetup{
    colorlinks=true,
    linkcolor=blue,
    urlcolor=blue,
    pdftitle={Relatorio do Dataset de Grupo - Analise de Dados e Modelacao em KNIME}
}

\newcommand{\placeholderfig}[2]{
\begin{figure}[H]
    \centering
    \fbox{%
        \begin{minipage}[c][6cm][c]{0.9\textwidth}
            \centering
            \textbf{[INSERIR IMAGEM KNIME]}
        \end{minipage}}
    \caption{#2}
\end{figure}
}

\title{\textbf{Relatorio de Resultados}\\
Dataset de Grupo: Spotify Tracks\\
\large Analise de Dados e Modelacao em KNIME}
\author{\textbf{Unidade Curricular:} Analise e Descoberta de Informacao \\
\textbf{Grupo:} [Preencher] \\
\textbf{Elementos:} [Preencher]}
\date{\today}

\begin{document}

\maketitle
\thispagestyle{empty}
\clearpage
\tableofcontents
\clearpage
\pagenumbering{arabic}

\section{Introducao}
Este relatorio documenta o trabalho realizado sobre o \textbf{dataset de grupo} do Spotify, com implementacao em \textit{workflows} KNIME e enquadramento completo no metodo CRISP-DM. O objetivo foi descrever um processo analitico completo, focado em duas tarefas distintas e complementares: \textbf{(i) clustering} para descobrir perfis acusticos e \textbf{(ii) regressao supervisionada} para prever \texttt{popularity}. A avaliacao foi feita no mesmo dataset, com particao treino/teste e interpretacao adicional dos clusters por genero.

O texto foi escrito para refletir fielmente o que existe no workflow, evitando passos que nao foram implementados.

\section{Dominio, objetivos e abordagem}
\subsection{Dominio analisado}
O dominio escolhido e o da analise musical e de consumo digital. As caracteristicas acusticas podem apoiar segmentacao editorial, recomendacao e analise de perfis sonoros. Variaveis como \texttt{danceability}, \texttt{energy}, \texttt{loudness}, \texttt{acousticness} e \texttt{tempo} capturam propriedades musicais diretas, enquanto \texttt{genre}, \texttt{release\_year} e \texttt{explicit} adicionam contexto de catalogo.

\subsection{Objetivos}
Os objetivos do trabalho sao:
\begin{itemize}[leftmargin=*]
    \item Garantir uma preparacao robusta dos dados (duplicados, missing, outliers e features derivadas);
    \item Descobrir perfis acusticos com \textit{k-means} e interpretar a composicao dos clusters;
    \item Treinar modelos de regressao para prever \texttt{popularity} com base nas features acusticas;
    \item Produzir uma analise critica dos resultados, focada na interpretacao dos clusters e no desempenho de regressao.
\end{itemize}

\subsection{Estrategia seguida}
A estrategia foi iterativa e alinhada com CRISP-DM: leitura e exploracao dos dados, preparacao sequencial, modelacao em ramos separados e avaliacao com foco na regressao de \texttt{popularity}. Esta abordagem evita comparacoes indevidas entre tarefas diferentes (clustering vs regressao) e reforca a interpretacao baseada em dados.

\section{Metodologia CRISP-DM}
\subsection{Compreensao do negocio}
O objetivo analitico e caracterizar musica por perfis acusticos e prever \texttt{popularity}. A utilidade pratica inclui segmentacao de catalogo e leitura critica do que os atributos sonoros explicam sobre popularidade.

\subsection{Compreensao dos dados}
O workflow confirma a estrutura do CSV, identifica duplicados, missing e valores fora do dominio e valida a distribuicao dos generos. A exploracao inclui estatisticas descritivas, histogramas, boxplots e distribuicao por genero.

\subsection{Preparacao dos dados}
Foi implementado um pipeline sequencial: remocao de duplicados, saneamento de valores invalidos, imputacao de missing, tratamento de outliers e criacao de features derivadas. O objetivo e garantir consistencia das features e reduzir ruido antes da modelacao.

\subsection{Modelacao}
Foram criados dois ramos principais:
\begin{itemize}[leftmargin=*]
    \item \textbf{Clustering}: \textit{k-means} com features acusticas normalizadas, avaliacao por elbow, silhouette e PCA.
    \item \textbf{Regressao supervisionada}: modelos linear e arvore de regressao para \texttt{popularity}, com normalizacao consistente e avaliacao com \textit{Numeric Scorer}.
\end{itemize}

\subsection{Avaliacao}
Avaliacao quantitativa: metricas de regressao (R2, MAE e RMSE) no teste. Avaliacao qualitativa: interpretacao de clusters por genero real, evitando comparacao direta entre tarefas diferentes.

\subsection{Deployment e uso}
Os artefactos resultantes (modelo de clustering, modelos de regressao e tabelas de interpretacao) podem ser reutilizados para analise de novos datasets com a mesma estrutura, mantendo o mesmo pipeline de preparacao e normalizacao.

\section{Descricao do dataset}
\subsection{Estrutura geral}
O ficheiro \texttt{spotify\_tracks.csv} contem \textbf{100\,500 registos} e \textbf{21 colunas}. O periodo temporal vai de \textbf{2000 a 2024} e o dataset cobre \textbf{20 generos} musicais. O ficheiro \texttt{spotify\_tracks\_b.csv} tem a mesma estrutura e esta disponivel para teste externo, mas nao foi usado no workflow final. Existe ainda o ficheiro \texttt{spotify\_tracks\_errors\_log.csv}, com o registo de erros introduzidos no gerador do dataset, usado para validar a preparacao de dados.

\begin{table}[H]
\centering
\caption{Estrutura do dataset de grupo}
\begin{tabular}{p{4cm}p{10cm}}
\toprule
Grupo de variaveis & Exemplos \\
\midrule
Identificacao & \texttt{track\_id}, \texttt{track\_name}, \texttt{artist\_name}, \texttt{album\_name} \\
Contexto & \texttt{release\_year}, \texttt{genre}, \texttt{explicit}, \texttt{key}, \texttt{mode}, \texttt{time\_signature} \\
Target & \texttt{popularity} \\
Audio & \texttt{danceability}, \texttt{energy}, \texttt{loudness}, \texttt{speechiness}, \texttt{acousticness}, \texttt{instrumentalness}, \texttt{liveness}, \texttt{valence}, \texttt{tempo}, \texttt{duration\_ms} \\
\bottomrule
\end{tabular}
\end{table}

\subsection{Qualidade dos dados}
O dataset nao foi criado para ser limpo; foi criado para simular um ambiente realista de preparacao. O gerador introduz valores em falta, outliers, violacoes de dominio e duplicados. Isto significa que a analise nao podia começar diretamente pela modelacao. Era necessario demonstrar, com evidencia, que os problemas estavam identificados e que o workflow sabia resolvê-los de forma consistente.

Em termos praticos, os problemas mais relevantes sao os seguintes. Existem valores em falta nas features de audio, o que afeta tanto a analise exploratoria como a modelacao. Existem linhas duplicadas completas, que distorcem contagens, distribuicoes e treino de modelos. Existem ainda valores fora do dominio esperado em variaveis limitadas a intervalos bem definidos, como as features normalizadas entre 0 e 1, e outliers fortes em variaveis como \texttt{tempo}, \texttt{loudness} e \texttt{duration\_ms}.

\subsection{Exploracao inicial}
A fase de exploracao foi organizada em torno de varios blocos visuais no KNIME. O objetivo era verificar rapidamente se o dataset estava lido corretamente, perceber a forma das distribuicoes e encontrar sinais de problemas a tratar.

\begin{figure}[H]
    \centering
    \includegraphics[width=0.9\textwidth]{imgs/datapopularidade.png}
    \caption{Exploracao univariada do dataset Spotify}
\end{figure}

\begin{figure}[H]
    \centering
    \includegraphics[width=0.9\textwidth]{imgs/loudeness.png}
    \caption{Outliers e dispersao inicial}
\end{figure}

\begin{figure}[H]
    \centering
    \includegraphics[width=0.9\textwidth]{imgs/dataplot.png}
    \caption{Relacoes bivariadas entre features acusticas}
\end{figure}

\begin{figure}[H]
    \centering
    \includegraphics[width=0.95\textwidth]{imgs/clusters.png}
    \caption{Representatividade por genero}
\end{figure}

A exploracao mostrou que as variaveis acusticas tinham escalas muito diferentes, o que e natural, mas tambem confirmou que havia varios pontos de anomalia a tratar. Por exemplo, alguns valores de \texttt{danceability}, \texttt{valence} ou \texttt{energy} surgem fora do intervalo esperado. Em paralelo, o comportamento de \texttt{tempo} e \texttt{loudness} inclui valores extremos que justificam tratamento especifico. Esta leitura inicial explica porque o workflow nao passa diretamente da leitura do CSV para a modelacao.

\section{Preparacao dos dados em KNIME}
A preparacao foi desenhada como um pipeline sequencial, onde cada no acrescenta uma operacao concreta ao estado dos dados.

\subsection{Remocao de duplicados}
O primeiro passo foi o \texttt{Duplicate Row Filter}. A escolha de tratar duplicados antes de qualquer outra transformacao evita que a limpeza ou a imputacao sejam aplicadas duas vezes aos mesmos registos. O workflow compara praticamente todas as colunas relevantes e conserva a primeira ocorrencia, mantendo a ordem original. O resultado observado e claro: de 100\,500 linhas iniciais, o dataset passa para 100\,000 linhas, o que confirma a existencia de 500 duplicados removidos no inicio do pipeline.

\subsection{Sanitizacao e engenharia de atributos}
Depois dos duplicados, o no \texttt{Expression} concentra varias operacoes que poderiam estar espalhadas por mais nodos. Aqui sao feitas tres coisas muito importantes. Primeiro, valores invalidos em variaveis limitadas a dominos conhecidos sao convertidos em missing para serem resolvidos mais tarde com regra consistente. Segundo, a variavel \texttt{explicit} e transformada para uma representacao numerica binaria. Terceiro, sao criadas features derivadas como \texttt{duration\_min} e a classe auxiliar \texttt{popularity\_class}.

Esta etapa e crucial porque transforma o dataset num formato mais adequado a aprendizagem automatica sem perder a informacao original de contexto. A criacao de \texttt{duration\_min} torna a duracao mais interpretabvel do que \texttt{duration\_ms}, e a variavel \texttt{popularity\_class} permite estratificar o split de treino/teste e apoiar a leitura dos resultados.

\subsection{Tratamento de valores em falta}
O no \texttt{Missing Value} recebe tanto os missing originais do CSV como os que sao gerados pelo passo anterior. A estrategia nao e uniforme; e precisamente essa diferenca que mostra maturidade metodologica. Para algumas features, o workflow usa media; para outras, mediana; e para o target \texttt{popularity}, a linha e removida quando o valor esta em falta.

\begin{table}[H]
\centering
\caption{Estrategias de imputacao no no Missing Value}
\begin{tabular}{p{4cm}p{8cm}}
\toprule
Coluna & Estrategia \\
\midrule
\texttt{danceability} & Mean \\
\texttt{valence} & Mean \\
\texttt{energy}, \texttt{speechiness}, \texttt{acousticness}, \texttt{instrumentalness}, \\
\texttt{liveness}, \texttt{loudness}, \texttt{tempo}, \texttt{duration\_ms} & Median \\
\texttt{popularity} & Remove Row \\
\bottomrule
\end{tabular}
\end{table}

Depois desta fase, o dataset fica com 99\,353 linhas. O valor e importante porque confirma que a preparacao nao removeu tudo indiscriminadamente; removeu apenas o que era metodologicamente necessario remover.

\subsection{Tratamento de outliers}
Os outliers continuos sao tratados num no especifico, \texttt{Numeric Outliers}, que opera sobre \texttt{duration\_ms}, \texttt{loudness} e \texttt{tempo}. O criterio de deteccao e o IQR, com scalar 3.0, e a acao escolhida e substituir o valor fora de intervalo pelo limite admissivel mais proximo. Esta decisao e relevante porque preserva a linha e impede que observacoes extremas deformem a modelacao.

\subsection{Selecao de variaveis e separacao de ramos}
Depois da limpeza principal, o workflow organiza uma base comum com 15 colunas. A partir daqui o workflow divide-se em dois ramos:
\begin{itemize}[leftmargin=*]
    \item \textbf{Ramo de clustering}: usa 10 features acusticas, com normalizacao Min-Max para [0,1].
    \item \textbf{Ramo de regressao}: mantem \texttt{popularity} como target, aplica normalizacao Min-Max e faz split treino/teste estratificado por \texttt{popularity\_class}.
\end{itemize}

\begin{figure}[H]
    \centering
    \includegraphics[width=0.9\textwidth]{imgs/datapreparation.png}
    \caption{Pipeline de preparacao no KNIME}
\end{figure}

\section{Modelacao}
\subsection{Clustering}
O clustering procura agrupar faixas com perfis sonoros semelhantes. Para isso o workflow usa apenas as features acusticas normalizadas, o que e a opcao correta porque o objetivo nao e agrupar por ano, genero ou identificador, mas sim por assinatura sonora.

\subsubsection*{Metodo do cotovelo}
A escolha do numero de clusters foi analisada com varias execucoes de k-Means para valores de \textit{k} entre 3 e 9. O workflow nao usa uma funcao automatica pronta para o elbow; em vez disso, reconstrói manualmente a soma das distancias quadradas dentro de cada cluster e agrega os resultados num grafico de linha. Esse detalhe e importante porque mostra que a decisao de \textit{k} nao foi arbitraria: foi suportada por uma curva de WCSS produzida no proprio KNIME.

\begin{figure}[H]
    \centering
    \includegraphics[width=0.9\textwidth]{imgs/line_plot.png}
    \caption{Metodo do cotovelo}
\end{figure}

\subsubsection*{k-Means final}
O modelo final de clustering foi fixado em \textbf{5 clusters}, com inicializacao aleatoria, seed 1 e maximo de 99 iteracoes. O output do proprio no inclui a coluna de cluster, pelo que nao foi necessario um assigner separado. Esta decisao final e coerente com a leitura da curva do elbow e com a interpretacao posterior dos centroides.

\begin{figure}[H]
    \centering
    \includegraphics[width=0.9\textwidth]{imgs/image.png}
    \caption{Clusters projetados em duas dimensoes (PCA)}
\end{figure}

\begin{figure}[H]
    \centering
    \includegraphics[width=0.9\textwidth]{imgs/k_means_features.png}
    \caption{Centroides dos clusters por feature}
\end{figure}

\subsubsection*{Interpretacao dos clusters}
Para tornar os clusters interpretaveis, o workflow devolve os centroides a forma longa, cria graficos de barras e junta o genero musical para verificar composicao de cada grupo. Esta parte e essencial porque um cluster sem leitura semantica tem pouco valor analitico.

A qualidade do agrupamento foi analisada adicionalmente com silhouette. O metodo hierarquico sobre uma amostra obteve um valor medio global de aproximadamente \textbf{0.0732}, o que e baixo e sugere sobreposicao entre grupos. Na comparacao entre k = 4, 5 e 6, os valores globais observados foram aproximadamente \textbf{0.174}, \textbf{0.178} e \textbf{0.161}, respetivamente, o que suporta a escolha de \textbf{k = 5} como compromisso.

\begin{figure}[H]
    \centering
    \includegraphics[width=0.7\textwidth]{imgs/4clusters.png}
    \caption{Silhouette para k = 4}
\end{figure}

\begin{figure}[H]
    \centering
    \includegraphics[width=0.7\textwidth]{imgs/5clusters.png}
    \caption{Silhouette para k = 5}
\end{figure}

\begin{figure}[H]
    \centering
    \includegraphics[width=0.7\textwidth]{imgs/6clusters.png}
    \caption{Silhouette para k = 6}
\end{figure}

\begin{figure}[H]
    \centering
    \includegraphics[width=0.95\textwidth]{imgs/clusters.png}
    \caption{Distribuicao dos generos por cluster}
\end{figure}

\subsubsection*{Comparacao hierarquica}
O workflow inclui ainda um ramo de clustering hierarquico com distancia euclidiana e linkage complete. Este ramo nao substitui o k-Means final, mas serve para comparar a qualidade da segmentacao. A leitura critica deste valor de silhouette e importante: o objetivo nao e apenas escolher um algoritmo, mas perceber se a estrutura dos dados suporta clusters nitidos ou se existe uma sobreposicao natural entre estilos musicais.

\subsection{Regressao supervisionada de popularidade}
O objetivo supervisionado e prever \texttt{popularity}. O workflow inclui um ramo de regressao linear e um ramo de arvore de regressao, ambos avaliados com \textit{Numeric Scorer}. A normalizacao e feita no treino e aplicada ao teste para evitar fuga de informacao.

\begin{figure}[H]
    \centering
    \includegraphics[width=0.95\textwidth]{imgs/regretion.png}
    \caption{Ramo de regressao no KNIME}
\end{figure}

\begin{figure}[H]
    \centering
    \includegraphics[width=0.9\textwidth]{imgs/regressionlinear.png}
    \caption{Regressao linear: real vs previsto}
\end{figure}

\begin{figure}[H]
    \centering
    \includegraphics[width=0.9\textwidth]{imgs/regressionlinearh.png}
    \caption{Regressao linear: distribuicao dos residuos}
\end{figure}

\begin{figure}[H]
    \centering
    \includegraphics[width=0.6\textwidth]{imgs/regressionlinearscorer.png}
    \caption{Metricas da regressao linear}
\end{figure}

\begin{figure}[H]
    \centering
    \includegraphics[width=0.9\textwidth]{imgs/regressiontreeplot.png}
    \caption{Arvore de regressao: real vs previsto}
\end{figure}

\begin{figure}[H]
    \centering
    \includegraphics[width=0.9\textwidth]{imgs/regressiontreeh.png}
    \caption{Arvore de regressao: distribuicao dos residuos}
\end{figure}

\begin{figure}[H]
    \centering
    \includegraphics[width=0.6\textwidth]{imgs/simpleregressionscorer.png}
    \caption{Metricas da arvore de regressao}
\end{figure}

\begin{figure}[H]
    \centering
    \includegraphics[width=0.9\textwidth]{imgs/randonforestplot.png}
    \caption{Random forest: real vs previsto}
\end{figure}

\begin{figure}[H]
    \centering
    \includegraphics[width=0.9\textwidth]{imgs/randomforesth.png}
    \caption{Random forest: distribuicao dos residuos}
\end{figure}

\begin{figure}[H]
    \centering
    \includegraphics[width=0.6\textwidth]{imgs/randomforestscorer.png}
    \caption{Metricas da random forest}
\end{figure}

\section{Resultados e analise critica}
\subsection{Resumo dos resultados mais relevantes}
\begin{table}[H]
\centering
\caption{Resumo dos principais resultados observados}
\begin{tabular}{p{5cm}p{8.5cm}}
\toprule
Elemento & Resultado observado \\
\midrule
Tamanho inicial do dataset & 100\,500 registos e 21 colunas \\
Duplicados removidos & 500 linhas \\
Base comum apos preparacao & 99\,353 linhas e 15 colunas \\
Ramo de clustering & 99\,353 linhas e 10 features normalizadas \\
Split treino/teste & 80/20, estratificado por \texttt{popularity\_class}, seed 1 \\
Silhouette hierarquico & aproximadamente 0.0732 \\
Silhouettes adicionais (k = 4, 5, 6) & aproximadamente 0.174, 0.178 e 0.161 \\
Cluster final & \textit{k} = 5 \\
Regressao linear & R2 de 0.018 e RMSE de 17.679 \\
Arvore de regressao & R2 de -1.092 e RMSE de 25.803 \\
Random forest & R2 de 0.0 e RMSE de 17.842 \\
\bottomrule
\end{tabular}
\end{table}

\subsection{Leitura critica do clustering}
Os resultados de clustering mostram que existe estrutura nos dados, mas a separacao nao e perfeita. Um silhouette perto de 0.07 e um sinal de clusters relativamente sobrepostos, o que faz sentido num dataset musical: as faixas podem partilhar intensidades semelhantes de energia, valencia e dancabilidade, sem que isso se traduza em grupos radicalmente distintos. Em termos praticos, isto significa que o clustering e mais util como ferramenta de segmentacao e exploracao do que como rotulagem rigorosa e definitiva.

Os clusters nao sao perfeitamente separados no espaco geometrico, mas ainda assim mantem uma identidade suficiente para uso operacional como segmentacao exploratoria.

\subsection{Leitura critica da regressao de popularidade}
Os resultados de regressao mostram capacidade limitada de explicacao da variabilidade de \texttt{popularity}. O R2 baixo na regressao linear sugere que a relacao entre features acusticas e popularidade e fraca ou nao linear, enquanto a arvore de regressao apresenta desempenho inferior neste estado. Estes sinais sao consistentes com um target influenciado por fatores externos ao audio (marketing, contexto cultural e exposicao).

\subsection{Limitações do estudo}
A principal limitacao e que o dataset e sintético, ainda que inspirado em distribuicoes realistas. Isso e util para treino, mas nao substitui completamente um catalogo real com variacoes de metadata mais complexas. Outra limitacao e a ausencia, no estado atual do workflow, de uma tabela unica consolidada com metrias de regressao finalizadas. Por fim, o trabalho analisa associacoes e padroes, mas nao pode afirmar causalidade entre as features e a popularidade.
Outra limitacao e a comparacao entre tarefas distintas: clustering e regressao respondem a perguntas diferentes. Por isso, a comparacao deve ser qualitativa (interpretacao dos clusters) e nao quantitativa por uma unica metrica.

\section{Sugestoes e recomendacoes}
Recomenda-se testar diferentes valores de \textit{k} e seeds para avaliar estabilidade dos clusters e acrescentar uma analise de percentagens por cluster (normalizacao por cluster). No ramo de regressao, seria util consolidar as metricas numa tabela comparativa unica e testar tuning de hiperparametros para reduzir o erro. Estas extensoes nao sao obrigatorias para concluir o relatorio, mas ajudam a consolidar a interpretacao.

\section{Reprodutibilidade}
Para reproduzir o trabalho, basta seguir a sequencia geral do workflow. Primeiro, abrir o ficheiro do KNIME relativo ao dataset de grupo. Depois, executar o bloco de \textit{data understanding} para confirmar leitura, distribuicoes e duplicados. Em seguida, correr o pipeline de preparacao na ordem definida: remocao de duplicados, expression, missing value, numeric outliers e filtragem final. Depois, correr separadamente o ramo de clustering e o ramo de regressao, garantindo que a normalizacao e a particao sao executadas antes dos modelos.

Os ficheiros que sustentam a reproducao sao o CSV principal, o log de erros, o script de geracao e o workflow KNIME. No caso de querer repetir a analise de raiz, convem confirmar se os caminhos dos ficheiros estao corretos e se as ligacoes no KNIME apontam para a pasta do projeto.

\section{Conclusao}
O trabalho desenvolvido sobre o dataset de grupo mostra um processo CRISP-DM completo: compreensao do problema, exploracao e preparacao, modelacao e avaliacao. O clustering permite segmentar perfis acusticos e a regressao supervisionada mede o quanto as features explicam \texttt{popularity}. A interpretacao por clusters complementa a avaliacao quantitativa e evita comparacoes indevidas entre tarefas diferentes. O resultado e um workflow coerente, reproduzivel e alinhado com boas praticas de analise de dados.

\end{document}
